\documentclass[11pt,a4paper]{article}
\usepackage[T1]{fontenc}
\usepackage[utf8]{inputenc}
\usepackage{lmodern}
\usepackage{amsmath,amssymb,amsthm}
\usepackage[margin=27mm]{geometry}
\usepackage{booktabs,tabularx}
\usepackage[hidelinks]{hyperref}
\hypersetup{pdftitle={A Gap Laboratory: Oscillators, Circles and Winding Sectors},pdfauthor={navstokgap research working notes}}
\newtheorem{proposition}{Proposition}
\newcommand{\dd}{\mathrm{d}}
\newcommand{\spec}{\operatorname{spec}}
\newcommand{\R}{\mathbb R}
\newcommand{\Z}{\mathbb Z}
\title{A Gap Laboratory:\\Oscillators, Circles and Winding Sectors}
\author{navstokgap research working notes}
\date{9 September 2026 --- M03 draft with reviewed M07 extension}
\begin{document}
\maketitle
\begin{abstract}
A free classical particle on a circle has minimum action $mw^2L^2/(2T)$
in winding sector $w$. Comparing sector minima produces the positive action
gap $mL^2/(2T)$. The quantum particle on the same circle has energy gap
$2\pi^2\hbar^2/(mL^2)$, which closes as the circumference grows.
For the oscillator, the fixed-endpoint action Hessian has eigenvalues
$m[(n\pi/T)^2-\omega^2]$, while the quantum Hamiltonian has energy spacing
$\hbar\omega$. We derive the domains, bounds and limiting behaviour of these
objects. The comparison isolates normalization, confinement and topological
sector selection as mechanisms for positive lower bounds, and identifies a
precise route to an action-valued spectral spacing in the quantum oscillator.
The relativistic Kepler model adds an angular-action threshold $k/c$;
softening its singular core closes that threshold at fixed coupling.
\end{abstract}

\section{Three objects to compare}
Fix mass $m>0$, duration $T>0$, and, for quantum models, an input action
parameter $\hbar>0$. The oscillator has frequency $\omega>0$; the circle has
circumference $L>0$. We use $h=2\pi\hbar$.

A quadratic fluctuation form $Q$ has a normalized lower bound
$Q[\eta]\ge\mu\|\eta\|^2$. A Hamiltonian with ground energy $E_0$ and ground
projection $P_0$ has excitation gap
\[
 \Delta_H=\inf\spec\bigl(H|_{\operatorname{Ran}(I-P_0)}\bigr)-E_0.
\]
For a classical path class split into sectors $\mathcal P_w$, let
$s_w=\inf_{q\in\mathcal P_w}S[q]$. Comparing the $s_w$ defines a third object:
a gap between sector minima. The choice of object fixes the normalization,
units and admissible comparisons.

The Schr\"odinger constructions below use the standard Fourier and oscillator
realizations. Teschl's first edition supplies conventions and spectral background
\cite[\S7.2, \S8.3]{Teschl2009}. The derivations here retain physical constants
and give the particular bounds used by the research programme.

\section{Oscillator action: duration and normalization}
Consider real paths with $q(0)=q(T)=0$ and
\[
 S[q]=\frac m2\int_0^T(\dot q^2-\omega^2q^2)\dd t.
\]
The rest path $q_*=0$ is a classical solution for every $T$. Its exact action
difference is $S[\eta]-S[0]=Q_T[\eta]/2$, where
\[
 Q_T[\eta]=m\int_0^T(\dot\eta^2-\omega^2\eta^2)\dd t,
 \qquad \eta\in H^1_0(0,T).
\]
Complexification gives a sesquilinear form. Its self-adjoint operator on
$L^2(0,T)$ is:
\[
 J_T=-m\frac{\dd^2}{\dd t^2}-m\omega^2,
 \qquad D(J_T)=H^2(0,T)\cap H^1_0(0,T).
\]
This specializes the classical second-variation and Jacobi framework
\cite[\S8.2]{Cristoferi2016CalculusVariations}.

\begin{proposition}[Duration-dependent fluctuation spectrum]\label{prop:hessian}
An orthonormal eigenbasis is
\[
 e_n(t)=\sqrt{2/T}\sin(n\pi t/T),\qquad
 \lambda_n=m[(n\pi/T)^2-\omega^2],\quad n\ge1.
\]
The best $L^2$ lower-bound constant is
$\mu_T=m[(\pi/T)^2-\omega^2]$.
The form is strictly positive for $T<\pi/\omega$, is nonnegative with a
one-dimensional kernel at $T=\pi/\omega$, and has a negative direction for
$T>\pi/\omega$. At $T=k\pi/\omega$, $k\ge1$, its kernel is spanned by $e_k$.
For every $T>0$, positive action differences from $q_*=0$ accumulate at zero.
\end{proposition}
\begin{proof}
The sine basis diagonalizes the Dirichlet Laplacian. Its diagonal domain is
the stated Sobolev domain, and adding the bounded scalar $-m\omega^2$ shifts
each eigenvalue. For $\eta\in H^1_0$ with coefficients $b_n$,
\[
 Q_T[\eta]=\sum_{n\ge1}\lambda_n|b_n|^2
 \ge\lambda_1\sum_{n\ge1}|b_n|^2.
\]
Equality is attained on $e_1$, proving optimality. The signs and kernels follow
by comparing $n$ with $\omega T/\pi$. In particular, the number of negative
eigenvalues is $\lceil\omega T/\pi\rceil-1$.
Choose an integer $n>\omega T/\pi$. Then
\[
 S[A\sin(n\pi t/T)]-S[0]
 =\frac{mA^2T}{4}[(n\pi/T)^2-\omega^2]>0\quad(A\ne0).
\]
This tends to zero with $A$, and the paths converge to zero in $H^1$.
\end{proof}

At a conjugate duration $T=k\pi/\omega$, the zero mode is itself a variation
through classical solutions $A\sin(\omega t)$ with the prescribed endpoints.
After the first conjugate duration, the rest path is a saddle. These changes
concern a boundary-value variational problem at selected duration.

With $\|\eta\|_2^2=\int_0^T|\eta|^2\dd t$, the coefficient $\mu_T$ has units
mass/time$^2$, and $Q_T$ has action units. If instead one uses the time-average
norm $T^{-1}\int_0^T|\eta|^2\dd t$, the coefficient is $T\mu_T$.
Fixing an amplitude norm is therefore part of interpreting a fluctuation bound.

\section{Oscillator Hamiltonian and an action-valued spectrum}
The quantum oscillator acts on $L^2(\R,\dd x)$ through
\[
 H_\omega=-\frac{\hbar^2}{2m}\frac{\dd^2}{\dd x^2}
           +\frac{m\omega^2x^2}{2}.
\]
We take its self-adjoint closure from finite sums of Hermite eigenfunctions.
The proof below constructs that realization and specifies its full domain.

\begin{proposition}[Confinement and action spacing]\label{prop:oscillator}
For $m,\omega,\hbar>0$ the oscillator has simple eigenvalues
\[
 E_n=\hbar\omega(n+\tfrac12),\qquad n=0,1,2,\ldots,
 \qquad \Delta_\omega=\hbar\omega.
\]
The operator $\widehat I=H_\omega/\omega$ has action-valued spectrum
$\hbar(n+\tfrac12)$ and excitation spacing $\hbar$.
At fixed $\hbar$, the energy gap closes as $\omega\downarrow0$; at fixed
$\omega$, it closes as $\hbar\downarrow0$.
\end{proposition}
\begin{proof}
Set $\ell=\sqrt{\hbar/(m\omega)}$ and $u=x/\ell$. The unitary change
$U\psi(u)=\sqrt\ell\,\psi(\ell u)$ gives
\[
 UH_\omega U^{-1}=\frac{\hbar\omega}{2}
       \left(-\frac{\dd^2}{\dd u^2}+u^2\right)
       =\hbar\omega(B^*B+\tfrac12),
 \quad B=\frac{u+\partial_u}{\sqrt2}.
\]
On polynomial times Gaussian functions, $B^*=(u-\partial_u)/\sqrt2$ and
$[B,B^*]=1$. Starting from $\phi_0=\pi^{-1/4}e^{-u^2/2}$, define
\[
 \phi_n=(n!)^{-1/2}(B^*)^n\phi_0.
\]
Integration by parts, $B\phi_0=0$ and the commutator give orthonormality,
$B^*B\phi_n=n\phi_n$, and the stated eigenvalues.

For completeness, a function $f\in L^2(\R)$ orthogonal to all $\phi_n$ is
orthogonal to every polynomial times $e^{-u^2/2}$. The function
\[
 F(z)=\int_\R\overline{f(u)}e^{-u^2/2}e^{zu}\dd u
\]
is entire: Cauchy--Schwarz gives integrable bounds for every derivative,
uniformly on compact $z$ sets. All derivatives at zero vanish, so $F=0$.
On the imaginary axis this is the Fourier transform of the integrable function
$\overline f e^{-u^2/2}$; Fourier uniqueness gives $f=0$.

Put $\psi_n=U^{-1}\phi_n$. In this basis the full operator domain is
\[
 D(H_\omega)=\left\{\psi=\sum_{n\ge0}c_n\psi_n:
                \sum_{n\ge0}E_n^2|c_n|^2<\infty\right\},
 \qquad H_\omega\psi=\sum_{n\ge0}E_nc_n\psi_n.
\]
A real diagonal operator on this maximal domain is self-adjoint. Finite
partial sums converge in its graph norm, identifying the claimed closure.
The spectrum and both gap limits now follow from the diagonal entries.
\end{proof}

For any normalized state in the form domain, expansion in this basis gives
\begin{equation}\label{eq:osc-coercivity}
 \langle H_\omega\rangle_\psi-E_0
 \ge\hbar\omega(1-|\langle\psi_0,\psi\rangle|^2).
\end{equation}
The expectation denotes the quadratic-form value when the state is outside
the operator domain. It is a sharp bound: equality holds in the span of the
ground and first excited states. Orthogonality to the ground yields the full
gap; an arbitrarily small excited component gives a correspondingly small
positive mean energy excess.

The classical reduced action variable is $I=(2\pi)^{-1}\oint p\dd q=E/\omega$.
Indeed the energy ellipse
$p^2/(2m)+m\omega^2q^2/2=E$ has area $2\pi E/\omega$.
Its values vary continuously with $E\ge0$. The quantized operator above has
spacing $\hbar$, or spacing $h$ for $2\pi\widehat I$.
For one classical period $P=2\pi/\omega$, Hamilton's time-integrated action is
instead
\[
 S_{\rm period}=\oint p\dd q-EP=0.
\]
This identifies precisely which action-like object acquires the quantum spacing.
The inputs are the Schr\"odinger state space and the specified $\hbar$.

\section{Free motion: circle, interval and line}
For a circle of circumference $L$, use $L^2(0,L)$ and periodic conditions on
both the function and its derivative:
\[
 H_L=-\frac{\hbar^2}{2m}\partial_x^2,\quad
 D(H_L)=\{\psi\in H^2(0,L):\psi(L)=\psi(0),\ \psi'(L)=\psi'(0)\}.
\]
This selects the untwisted periodic sector.
The momentum and energy quantisation from periodicity is the standard ring
calculation \cite[Lecture 23, chapter 2]{YalePHYS201Lecture23}.

\begin{proposition}[Spatial-size gap and its closure]\label{prop:circle}
The circle has normalized eigenfunctions and energies
\[
 v_n(x)=L^{-1/2}e^{2\pi inx/L},\qquad
 E_n=\frac{2\pi^2\hbar^2n^2}{mL^2},\qquad n\in\Z.
\]
The ground state is the constant $v_0$, and its excitation gap is
\[
 \Delta_L=\frac{2\pi^2\hbar^2}{mL^2}.
\]
For the Dirichlet interval $(0,L)$ the energies are
$E_n^{\rm D}=\pi^2\hbar^2n^2/(2mL^2)$, $n\ge1$, and the ground-to-first-excited
gap is $3\pi^2\hbar^2/(2mL^2)$. Both gaps tend to zero as $L\to\infty$.
On the line, $H_\R=-\hbar^2\partial_x^2/(2m)$ with domain $H^2(\R)$ has
spectrum $[0,\infty)$ and an empty point spectrum.
\end{proposition}
\begin{proof}
Fourier series diagonalize the circle operator. The domain condition is
\[\sum_{n\in\Z}n^4|c_n|^2<\infty,\]
which defines the periodic $H^2$ domain.
The complete exponential basis proves the spectrum and self-adjointness.
The first excitation consists of the $n=\pm1$ pair. Replacing exponentials
by the complete sine basis and using domain $H^2(0,L)\cap H^1_0(0,L)$ proves
the interval statements; its excitation gap is $E_2^{\rm D}-E_1^{\rm D}$.

On the line the unitary Fourier transform converts the operator into
multiplication by $\hbar^2k^2/(2m)$ on its maximal domain. The essential range
is $[0,\infty)$, and every level set has Lebesgue measure zero, so there are
no square-integrable eigenvectors. For any $\epsilon>0$, a normalized Fourier
function supported in
$0<|k|<\sqrt{2m\epsilon}/\hbar$ has energy expectation in $(0,\epsilon)$.
\end{proof}

In particular the line's zero spectral edge has a different ground-state
structure from the circle's normalized constant ground state. The comparison
of gaps above gives their exact dependence on spatial size; an operator limit
between varying Hilbert spaces also requires a specified identification.

For a circle form-domain vector, Parseval gives the sharp estimate
\begin{equation}\label{eq:circle-coercivity}
 \frac{\hbar^2}{2m}\int_0^L|\psi'|^2\dd x
 \ge\Delta_L\left(\|\psi\|^2-|\langle v_0,\psi\rangle|^2\right).
\end{equation}
The same periodic Poincar\'e constant $\kappa_L=(2\pi/L)^2$ gives mean-zero
heat relaxation rate $\nu\kappa_L$ and quantum energy gap
$\hbar^2\kappa_L/(2m)$. Their units and generators fix their different roles.
The $L^{-2}$ factor closes both bounds in growing volume.

\section{A positive classical action gap between winding sectors}
Describe a circle-valued loop based at zero by its absolutely continuous lift
$q\in H^1(0,T;\R)$, with $q(0)=0$ and $q(T)=wL$ for $w\in\Z$.
This makes the winding sectors explicit:
\[
 \mathcal P_w=\{q\in H^1(0,T;\R):q(0)=0,\ q(T)=wL\},
 \quad S[q]=\frac m2\int_0^T\dot q^2\dd t.
\]

\begin{proposition}[Sector minimum action]\label{prop:winding}
The unique minimizing lift in sector $w$ is $q_w(t)=wLt/T$, with
\[
 s_w=S[q_w]=\frac{mw^2L^2}{2T},\qquad
 g_{\rm wind}:=\inf_{w\ne0}(s_w-s_0)=\frac{mL^2}{2T}>0.
\]
These minima are also the action values of all free classical solutions with
the prescribed circle endpoints and duration. Within each individual sector,
positive action excesses above its minimum accumulate at zero.
\end{proposition}
\begin{proof}
Cauchy--Schwarz yields
\[
 (wL)^2=\left(\int_0^T\dot q\dd t\right)^2
 \le T\int_0^T\dot q^2\dd t.
\]
Equality holds exactly when $\dot q$ is constant almost everywhere, fixing
the minimizing lift. Free classical solutions solve $\ddot q=0$, hence are
precisely these constant-speed lifts. Minimizing $w^2$ over nonzero integers
gives the gap. For $\eta\in H^1_0(0,T)$ the cross term integrates to zero:
\[
 S[q_w+\eta]-S[q_w]=\frac m2\int_0^T\dot\eta^2\dd t.
\]
Taking $\eta=A\sin(\pi t/T)$ gives $mA^2\pi^2/(4T)$, positive for $A\ne0$
and approaching zero with $A$.
\end{proof}

The positive bound concerns nonzero winding, or comparison of the stationary
solutions across sectors. It depends on the global length $L$ and duration $T$:
$g_{\rm wind}\to0$ as $T\to\infty$ at fixed $L$, or as $L\downarrow0$ at fixed
$T$. At fixed $T$ it grows with $L$, while the quantum circle gap decreases.
Their exact product is
\[
 g_{\rm wind}\Delta_L=\frac{\pi^2\hbar^2}{T}.
\]
The product records reciprocal $L$ scaling for two different quantities.

\section{Finite speed in the winding example}
Finite $c>0$ can be kept as an independent input. For the standard relativistic
free-particle Lagrangian \cite{FeynmanLectures}, measure action relative to rest:
\[
 \Delta S_c[q]=mc^2\int_0^T\left(1-\sqrt{1-\dot q^2/c^2}\right)\dd t.
\]
Use absolutely continuous lifted paths with $|\dot q|<c$ almost everywhere.
A winding sector is admissible exactly when $|w|L<cT$. Necessity follows
from $|w|L\le\int_0^T|\dot q|\dd t<cT$; a constant-speed lift proves
sufficiency. The velocity integrand is strictly convex, with second derivative
$m(1-v^2/c^2)^{-3/2}$. Jensen's inequality therefore gives its unique minimum:
\[
 s^c_w=mc^2T\left(1-\sqrt{1-\frac{w^2L^2}{c^2T^2}}\right).
\]
When $L<cT$, the first nonzero sector has gap
\[
 g^c_{\rm wind}=mc^2T\left(1-\sqrt{1-\frac{L^2}{c^2T^2}}\right)
 =\frac{mL^2}{2T}+\frac{mL^4}{8c^2T^3}+O(c^{-4})
\]
at fixed $m,L,T$ as $c\to\infty$. For $L\ge cT$ the admissible class contains
only the zero-winding sector. This is absence of a comparison sector, so the
finite gap formula applies precisely on $L<cT$.
At fixed finite $c,L$, increasing $T$ again closes the gap.

\section{A singular-core angular-action threshold}
The relativistic Kepler model has a positive infimum of angular action among
regular bound trajectories. Softening its central singularity closes that
infimum at fixed coupling and speed limit. For this section write angular
momentum as $J$ to reserve $L$ for the circle's circumference. On the planar
canonical domain $r>0$, use
\[
 H=\sqrt{m^2c^4+c^2(p_r^2+J^2/r^2)}-k/r,\qquad m,k,c>0.
\]
The potential is static and external; recoil, radiation and field dynamics
are absent. Energy includes rest energy. A regular bound trajectory is complete
and has radius in a compact subset of $(0,\infty)$. Normalize the angular
canonical cycle by $I_\theta=(2\pi)^{-1}|\int_0^{2\pi}J\dd\theta|=|J|$,
which does not require a closed spatial rosette. Define
\[
 g_\theta=\inf\{|J|:\text{a regular bound trajectory exists}\}.
\]
This is an admissibility threshold with action units, since $[k/c]$ is action.

\begin{proposition}[Kepler threshold and core dependence]\label{prop:kepler}
For the singular potential, regular bound motion exists exactly for
\[
 |J|>k/c,\qquad
 mc^2\sqrt{1-k^2/(c^2J^2)}\le E<mc^2.
\]
The lower energy endpoint is circular. Thus $g_\theta=k/c$ is an excluded
infimum, and the admitted angular actions above it are continuous.
For $V_a(r)=-k/\sqrt{r^2+a^2}$ with any fixed $a>0$, every $|J|>0$
admits a regular bound circle, so $g_\theta(a)=0$.
\end{proposition}
\begin{proof}
Put $\ell=|J|$ and $x=1/r$. At zero radial momentum the singular energy is
\[
 F_\ell(x)=\sqrt{m^2c^4+c^2\ell^2x^2}-kx.
\]
For $\ell>0$ it is strictly convex, with derivative $-k$ at zero and limiting
derivative $c\ell-k$ at infinity. A finite minimum therefore exists exactly
when $c\ell>k$. Solving $F'_\ell=0$ gives the stated energy and radius
$r_0=\ell\sqrt{c^2\ell^2-k^2}/(mkc)$.
For energies strictly between this minimum and $mc^2$, the allowed radial
interval lies between two finite positive turning radii. Energy also bounds
the momentum, so smooth Hamiltonian continuation gives completeness.
The minimum gives a circle; $E\ge mc^2$ has no finite outer turning radius.

At $c\ell=k$, $F_\ell$ decreases to zero at the excluded centre; below
threshold it decreases to minus infinity, with the $\ell=0$ case immediate.
There is no well. The inward branch satisfies
\[
 c^2p_r^2=E^2-m^2c^4+2Ek/r+(k^2-c^2\ell^2)/r^2,
 \qquad \dot r=c^2p_r/(E+k/r).
\]
Below threshold, its speed tends to $c\sqrt{k^2-c^2\ell^2}/k$.
At threshold allowed energies have $E>0$, and
$|\dot r|\sim c\sqrt{2E/k}\sqrt r$. Both yield finite collision time;
no collision continuation belongs to the specified regular domain.
This recovers Boyer's classification
\cite[Eqs.~(19)--(20), (37)--(42)]{Boyer2004UnfamiliarTrajectories}.

For the softened core and any $\ell>0$, the effective energy is
\[
 U_{a,\ell}(r)=\sqrt{m^2c^4+c^2\ell^2/r^2}-k/\sqrt{r^2+a^2}.
\]
It diverges at zero and has large-radius expansion
$mc^2-k/r+\ell^2/(2mr^2)+O(r^{-3})$.
It thus takes values below its limit $mc^2$ and attains a global minimum
at finite positive radius. With $p_r=0$, the vanishing radial derivative
gives a complete circular Hamiltonian trajectory. All positive $\ell$ are
therefore admitted. For any energy between the minimum and $mc^2$, allowed
radial components are compact as well.
\end{proof}

At fixed $k$, the singular threshold closes as $c\to\infty$; at fixed $c$
it closes as $k\downarrow0$. The latter is a limit through positive couplings:
the zero-coupling model has no nonstationary bound orbits. At fixed $\ell,k,m$, the circular
radius tends to $\ell^2/(mk)$ and the rest-subtracted minimum energy tends to
$-mk^2/(2\ell^2)$ in the Newtonian limit.

The core limit is singular even at fixed coupling. For each fixed
$0<\ell<k/c$, choose $y$ large enough that
$c\ell/y<k/\sqrt{1+y^2}$. Then
\[
 \lim_{a\downarrow0}aU_{a,\ell}(ay)
 =c\ell/y-k/\sqrt{1+y^2}<0.
\]
The minimum tends to minus infinity, while on $r\ge\delta>0$ the energy is
bounded below by $mc^2-k/\delta$. Hence every minimizing radius tends to zero.
These circles leave the singular model's regular domain. Consequently
$\lim_{a\downarrow0}g_\theta(a)=0$, whereas $g_\theta=k/c$ at the singular
potential. This limit is pointwise in fixed subcritical $\ell$.

\section{A bound-orbit action estimator}
Uniform circular phase gives a canonical covariance area equal to orbital
action, while displacement transport vanishes at both window extremes
(C054--C055). Write $q=R(\cos\theta,\sin\theta)$,
$p=P(-\sin\theta,\cos\theta)$, $\theta=\omega t+\Phi$, with uniform
$\Phi$ and $\ell=RP=\gamma mR^2\omega$,
$\gamma=(1-R^2\omega^2/c^2)^{-1/2}$.
For $x=q_1$, phase averaging gives
$\mathbb E[x(t)x(t+\Delta)]=(R^2/2)\cos z$, $z=\omega\Delta$, hence
\[
 \mathsf h_x(\Delta)=\frac m\Delta\operatorname{Var}(x(t+\Delta)-x(t))
 =\frac{\ell}{\gamma}\frac{1-\cos z}{z}.
\]
Given $x(t)=x$, the equal-weight next positions are
$x\cos z\pm\sqrt{R^2-x^2}\sin z$. The conditional variance is
$(R^2-x^2)\sin^2z$, with continuous zero endpoint extension; its
phase-averaged action coefficient is $\ell\sin^2z/(2\gamma z)$.
Both coefficients vanish at zero and infinite window. Since
$1-\cos z\le\min(z^2/2,2)$,
\[
 \mathsf h_x\ge\eta\ell\quad\Longrightarrow\quad
 2\gamma\eta\le z\le2/(\gamma\eta),\qquad\eta>0.
\]
Every interval above that fixed positive fraction has endpoint ratio at most
$1/(\gamma^2\eta^2)$; a broad plateau needs additional time scales.

For centered covariance $\Sigma$ of $(x,p_x)$, define
$\mathcal A_{\rm cov}=2\sqrt{\det\Sigma}$. Uniform phase gives
$\Sigma=\operatorname{diag}(R^2/2,P^2/2)$, hence
$\mathcal A_{\rm cov}=RP=\ell$. The projected ellipse area is $\pi RP$,
half the full planar canonical orbit integral. This fixes the factor two
relative to rms emittance\cite{Groening2018Busch}; here momenta have physical
canonical units, while the cited beam convention normalizes momenta.

An affine canonical map in this plane has $\Sigma'=S\Sigma S^T$ and
$\det S=1$, preserving the determinant. The projected orbit map is
\[
 S_t=\begin{pmatrix}\cos\omega t&(R/P)\sin\omega t\\
 -(P/R)\sin\omega t&\cos\omega t\end{pmatrix}.
\]
It preserves covariance area for every phase law, including the zero area
of a point phase with $\ell>0$. Uniform phase, or matching first and second
moments, suffices for equality to orbital action. This is conservation of a
prepared value on the fixed orbit, rather than attraction to a common value.
Uniform singular-Kepler circles have $\inf\mathcal A_{\rm cov}=k/c$;
any fixed softened core has infimum zero by the preceding proposition.
The estimator transfers the threshold with its preparation and core premises.

\section{Dilation closure of the model class}
A model class admitting arbitrarily small simultaneous space/time contractions
has zero infimum for a positive finite degree-one action observable
(C056--C057). Fix $H(q,p)=T(p)+V(q)$, where $T$ is either Newtonian or
relativistic kinetic energy, and $V\in C^2(D)$. For $a>0$ set
\[
 V_a(Q)=V(Q/a),\qquad q_a(t)=a q(t/a),\qquad p_a(t)=p(t/a).
\]
Differentiation gives $\dot q_a=\nabla T(p_a)$ and
$\dot p_a=-\nabla V_a(q_a)$ on $aD$. Mass, velocity and energy values
are preserved, while force and acceleration scale by $a^{-1}$.
The map $F_a(q,p)=(aq,p)$ multiplies the canonical two-form by $a$;
it is conformally symplectic. An invariant preparation pushes forward to an
invariant preparation because $\phi_a^tF_a=F_a\phi^{t/a}$.

Changing variables in the orbit and finite-time integrals yields
$J_a=aJ$ and $S_a=aS$. For the centered covariance in the fixed canonical
plane, $\Sigma_a=D_a\Sigma D_a$ with $D_a=\operatorname{diag}(a,1)$, hence
\[
 \mathcal A_{{\rm cov},a}=a\mathcal A_{\rm cov},\qquad
 \mathsf h_a(\Delta)=a\mathsf h(\Delta/a).
\]
Finite second moments are assumed. A finite long-window plateau therefore
scales by $a$, taking the window limit at fixed $a$ first.
All these actions have units of mass times length squared per time.

If the class contains one $0<A(M)<\infty$ and its contractions $M_a$ for
arbitrarily small $a$, then the positive values $aA(M)$ approach zero.
If a common finite value $K$ is required across such a class, $K=aK$
forces $K=0$. No regular model at $a=0$ is needed.
The infimum is over model/preparation pairs, not an operator spectrum.

For a Kepler potential this map sends $k$ to $ak$; for a softened core it
sends $(k,b)$ to $(ak,ab)$; for a harmonic potential it sends the spring
constant to $k_s/a^2$. Thus fixed coupling or a uniform force ceiling can
exclude this map. A particle speed ceiling and uniform circular phase
preserve it. The next test holds the potential and force ceiling fixed and
examines small-radius circles. Breaking dilation closure alone does not
prove a positive action floor. The external-potential construction imposes
no propagation law for a mediator.

\section{Fixed force ceiling and circular excitation}
A single smooth confining potential with a global force ceiling still admits
positive circular actions tending to zero (C058). Fix $m,c,K,b>0$ and take
\[
 V(r)=Kb^2(\sqrt{1+r^2/b^2}-1),\qquad
 f(r)=V'(r)=\frac{Kr}{\sqrt{1+r^2/b^2}}<Kb.
\]
On Cartesian configuration space the potential is smooth at the origin and
confining. Finite energy bounds position and momentum, so trajectories are
complete. Choose $Kb\le F_{\max}$ for a global force ceiling.
Relativistic circular balance is $Pv/R=f(R)$, the established central-force
relation\cite{Aguirregabiria2004Circular}. Set $s=Rf(R)$ and solve it as
\[
 \gamma_R=\tfrac12\left(s/(mc^2)+\sqrt{[s/(mc^2)]^2+4}\right),
 \quad v_R=c\sqrt{1-\gamma_R^{-2}},\quad \ell_R=mR\gamma_Rv_R.
\]
These give an exact circle for every $R>0$. At fixed $\ell_R$, differentiating
the radial energy gives
\[
 U''_{\ell_R}(R)=f'(R)+(3-v_R^2/c^2)f(R)/R>0,
 \qquad f'(R)=K(1+R^2/b^2)^{-3/2}.
\]
Thus each circle is a strict radial minimum, stable under small radial
perturbations at fixed angular momentum. For fixed parameters, expansion at
$R\downarrow0$ yields
\[
 v_R\sim\sqrt{K/m}R,\qquad \ell_R\sim\sqrt{mK}R^2,\qquad
 \omega_R\to\sqrt{K/m},\qquad H-mc^2\sim KR^2.
\]
Force and coordinate acceleration tend to zero; the period tends to a finite
constant. All sufficiently small circles meet any positive speed and
acceleration ceilings. Uniform phase transfers the zero action infimum to
$\mathcal A_{\rm cov}$. The limit is the smooth central equilibrium.

An independent lower excitation condition supplies a positive bound (C059).
For any relativistic circle with $f(R)\le F_{\max}$ and $v\ge v_*>0$,
$v_*<c$, balance gives
\[
 \ell=\frac{P^2v}{f(R)}
 \ge\frac{m^2v_*^3}{F_{\max}(1-v_*^2/c^2)}>0.
\]
Here $m^2v^3/(1-v^2/c^2)$ is increasing on $(0,c)$. The units are action;
the Newtonian formula omits the relativistic denominator. The floor is a
separate kinetic-excitation premise. A18 will test a corresponding bound
for general closed regular trajectories using their total turning.

\section{What the laboratory selects next}
\begin{center}
\begin{tabularx}{\linewidth}{@{}lXl@{}}
\toprule
Object & Lower bound or spacing & Gap-closing parameter \\
\midrule
Oscillator Hessian & $m[(\pi/T)^2-\omega^2]$ in the $L^2$ norm & $T\uparrow\pi/\omega$ \\
Quantum oscillator & $\Delta_\omega=\hbar\omega$ & $\omega\downarrow0$ \\
Oscillator action & spacing of $H_\omega/\omega$: $\hbar$ & $\hbar\downarrow0$ \\
Quantum circle & $\Delta_L=2\pi^2\hbar^2/(mL^2)$ & $L\to\infty$ \\
Classical winding & $g_{\rm wind}=mL^2/(2T)$ & $T\to\infty$ \\
Kepler angular action & $g_\theta=k/c$ & $c\to\infty$ \\
\bottomrule
\end{tabularx}
\end{center}

The oscillator supplies a confinement-generated energy gap in unbounded
configuration space. The circle supplies a geometry-generated quantum gap and
a classical action gap after a specified sector selection. Equations
\eqref{eq:osc-coercivity} and \eqref{eq:circle-coercivity} identify the exact
orthogonality condition that turns a spectral spacing into a statewise bound.

For the programme's next gap model, choose the sector, state space or coupling,
then ask whether the lower bound survives removal of the chosen duration or
size constraint. An interacting oscillator chain is a concrete next test of
uniformity in system size. The periodic Poincar\'e comparison also supplies a
precise companion to the fluid discussion: the shared estimate and its loss
at large volume can be followed explicitly.

The quantum spectra use a prescribed $\hbar$ and the stated Schr\"odinger
realizations. The classical sector bounds use topology and endpoint data.
Reconstructing a universal action parameter remains the separate physical-premise
question. The present results are standard solvable-model derivations assembled
to distinguish and test gap mechanisms.

The repository's \texttt{make check} validates source integrity, links and
citation metadata; written proofs supply mathematical verification.
\texttt{make papers} builds the manuscripts and checks references
and layout. The source companion records edition and reading coverage.
\bibliographystyle{plain}
\bibliography{references/library}
\end{document}
